\section{Contexto del proyecto y datos}

El backend computacional de este informe es el repositorio en \Julia{} utilizado para el flujo completo de EEG, desde la ingestión compatible con BIDS hasta la inferencia de conectividad. El proyecto separa configuración, módulos fuente, scripts de orquestación, datos de entrada/intermedios y outputs generados (\texttt{config/}, \texttt{src/}, \texttt{script/}, \texttt{data/}, \texttt{results/}). Esta separación es importante para el trabajo interno porque permite relanzamientos deterministas, cambios de parámetros controlados y una depuración más sencilla cuando una etapa concreta se comporta de forma inesperada.

En términos de datos, el análisis utiliza EEG en reposo de controles sanos y pacientes con esclerosis múltiple (EM), adquirido en condiciones EO y EC \citep{Mestre2021}. Operativamente, \textbf{EC} (eyes closed) corresponde a vigilia relajada con ojos cerrados, mientras que \textbf{EO} (eyes open) corresponde a reposo vigil con ojos abiertos y fijación visual. Las señales se representan como matrices canal-por-muestra y se propagan por cada fase con metadatos explícitos, de modo que la identidad del sujeto, la condición y la procedencia de cada run se conservan en todos los artefactos intermedios, de acuerdo con recomendaciones estándar de procesamiento EEG \citep{Yamada2018,Sanei2022,Cohen2014}.
\begin{figure}[H]
\centering
\includegraphics[width=0.58\textwidth]{methodology/fase0_resumen_sujetos.png}
\caption{Resumen operativo del estudio BRAIN usado en Fase 0: pacientes EM en T1/T2, controles y pérdidas de seguimiento.}
\label{fig:margin-fase0-summary}
\end{figure}

\begin{parametros}
\textbf{Parámetros canónicos de adquisición:} $C = 32$ canales, $F_s = \SI{500}{Hz}$ y $T = \SI{100}{s}$ por condición.
\end{parametros}
Para cada sujeto $s$ y condición $k \in \{\mathrm{EO},\mathrm{EC}\}$, el registro se representa como
\[
X_{s,k} \in \mathbb{R}^{C \times N},
\]
con $C$ canales y $N = F_s \cdot T$ muestras temporales por canal.
\begin{figure}[H]
\centering
\begin{tikzpicture}[x=0.72cm,y=0.58cm,>=Latex]
  \draw[rounded corners=2pt, draw=black!75, fill=black!3] (0,0) rectangle (5,3);
  \draw[black!30] (1,0)--(1,3); \draw[black!30] (2,0)--(2,3); \draw[black!30] (3,0)--(3,3); \draw[black!30] (4,0)--(4,3);
  \draw[black!30] (0,1)--(5,1); \draw[black!30] (0,2)--(5,2);
  \node at (2.5,1.5) {$X_{s,k}$};
  \node at (0.5,2.65) {$x_{1,1}$};
  \node at (4.5,2.65) {$x_{1,N}$};
  \node at (0.5,0.35) {$x_{C,1}$};
  \node at (4.5,0.35) {$x_{C,N}$};
  \draw[-Latex, thick] (0,-0.55)--(5,-0.55);
  \node[below] at (2.5,-0.55) {$N$ muestras};
  \draw[-Latex, thick] (-0.75,0)--(-0.75,3);
  \node[anchor=east] at (-0.82,1.5) {$C$ canales};
\end{tikzpicture}
\caption{Bloque matricial $X_{s,k}$: $s$ identifica sujeto y $k$ la condición EO/EC.}
\label{fig:data-matrix-block}
\end{figure}

La organización del dataset sigue convenciones BIDS, incluyendo jerarquía sujeto/sesión y metadatos EEG tipo sidecar \citep{Yamada2018}. Operativamente, las comprobaciones más críticas son la consistencia de etiquetas de canal entre ficheros, el etiquetado EO/EC de los runs y las referencias internas de BrainVision tras el renombrado (\texttt{DataFile=} y \texttt{MarkerFile=}). En la práctica, el preprocesamiento solo comienza cuando pasan cuatro controles: consistencia de nombres de canal, etiquetas EO/EC válidas, enlaces \texttt{.vhdr/.vmrk} actualizados y metadatos completos de participantes.

\begin{controlqc}
Antes de iniciar el preprocesamiento, cada run debe resolver correctamente el tríptico BrainVision (\texttt{.vhdr}, \texttt{.vmrk}, \texttt{.eeg}), tener condición EO/EC inequívoca, conservar etiquetas de canal compatibles con \texttt{electrodes.tsv} y pasar validación BIDS.
\end{controlqc}

\begin{fullwidth}
\begin{figure}[H]
\centering
\resizebox{0.93\fulltextwidth}{!}{\NSBIDSBrainVision}
\caption{Conversión conceptual desde archivos BrainVision hacia una estructura BIDS-EEG validable. El punto crítico no es solo renombrar, sino mantener enlaces internos coherentes, metadatos tabulares y validación reproducible.}
\label{fig:generated-bids-brainvision-flow}
\end{figure}
\end{fullwidth}

El tratamiento frecuencial se mantiene fijado a las mismas bandas utilizadas en el flujo original, lo que preserva compatibilidad con los módulos FFT, CSD y wPLI existentes y con outputs generados previamente \citep{Cohen2014,Sanei2022,Ruth2020}.

\begin{table}[H]
\centering
\caption{Bandas de frecuencia usadas en análisis espectral y conectividad.}
\label{tab:frequency-bands}
\begin{tabular}{@{}clc@{}}
\toprule
\multicolumn{2}{c}{\textbf{Banda}} & \textbf{Rango (Hz)} \\
\midrule
$\delta$ & Delta & 0.5--4.0 \\
$\theta$ & Theta & 4.0--8.0 \\
$\alpha$ & Alpha & 7.8--11.7 \\
$\beta_l$ & B. low & 12.0--15.0 \\
$\beta_m$ & B. mid & 15.0--18.0 \\
$\beta_h$ & B. high & 18.0--30.0 \\
$\gamma$ & Gamma & 30.0--50.0 \\
\bottomrule
\end{tabular}
\end{table}

\subsection{Disposición de electrodos y coordenadas}

Las posiciones de electrodos siguen el sistema internacional 10--20 \citep{Yamada2018,IriarteFranco2013}. La disposición usada en este proyecto incluye 31 canales EEG más referencia (FCz) y tierra (Fpz). Para cada sujeto/sesión, BIDS proporciona un fichero \bids{electrodes.tsv} (por ejemplo, \bids{sub-M05\_ses-T2\_electrodes.tsv}) con los campos \texttt{name}, \texttt{x}, \texttt{y}, \texttt{z} y \texttt{type}. Las coordenadas se normalizan y se usan posteriormente en visualización topográfica y en etapas relacionadas con CSD/fuente \citep{Cohen2014}.

En esta nomenclatura, la letra indica la región cortical aproximada: \textbf{F} (frontal), \textbf{C} (central), \textbf{T} (temporal), \textbf{P} (parietal) y \textbf{O} (occipital). Los números impares corresponden al hemisferio izquierdo, los pares al derecho, y \texttt{z} indica la línea media (por ejemplo, Fz, Cz, Pz). El sistema 10--20 implica que las posiciones adyacentes se definen mediante distancias relativas del 10\% o 20\% de medidas craneales estándar, lo que proporciona un muestreo espacial reproducible entre sujetos \citep{Yamada2018,IriarteFranco2013}.

Respecto a las etiquetas de tipo de canal, \texttt{EEG} denota electrodos activos de registro, \texttt{REF} denota el canal de referencia (aquí FCz) y \texttt{GND} denota el canal de tierra (aquí Fpz). Operativamente, \texttt{REF} es la referencia de voltaje usada por el amplificador para diferencias de potencial, mientras que \texttt{GND} estabiliza el circuito de adquisición y ayuda a reducir ruido eléctrico común \citep{Yamada2018}.
La Tabla~\ref{tab:electrode-layout-full} recoge la disposición completa de electrodos usada en el pipeline.

\begin{fullwidth}
\begin{figure}[H]
\centering
\resizebox{0.82\fulltextwidth}{!}{\NSEEGMontage}
\caption{Montaje EEG 10--20 usado como referencia visual del pipeline: 31 canales EEG, referencia FCz y tierra Fpz. El código de color agrupa regiones para facilitar revisión de coordenadas, interpretación de CSD y lectura de matrices de conectividad.}
\label{fig:generated-eeg-montage}
\end{figure}
\end{fullwidth}

\begin{table}[htbp]
\centering
\scriptsize
\caption{Nombres de electrodos, coordenadas normalizadas $(x,y,z)$ y tipo, ordenados por región.}
\label{tab:electrode-layout-full}
\begin{tabular}{@{}cc c c l@{}}
\toprule
\textbf{Nombre} & $x$ & $y$ & $z$ & \textbf{Tipo} \\
\midrule
\multicolumn{5}{@{}l}{\textit{Frontal}} \\
Fp2  & 0.35 &  0.90 & 0.260 & EEG \\
Fz   & 0.00 &  0.70 & 0.714 & EEG \\
F3   & -0.35 &  0.68 & 0.644 & EEG \\
F7   & -0.85 &  0.68 & 0.000 & EEG \\
F4   & 0.35 &  0.68 & 0.644 & EEG \\
F8   & 0.85 &  0.68 & 0.000 & EEG \\
FC5  & -0.60 &  0.68 & 0.421 & EEG \\
FC1  & -0.60 &  0.68 & 0.421 & EEG \\
FC6  & 0.60 &  0.68 & 0.421 & EEG \\
FC2  & 0.60 &  0.68 & 0.421 & EEG \\
FT9  & -0.95 &  0.20 & 0.240 & EEG \\
FT10 & 0.95 &  0.20 & 0.240 & EEG \\
\addlinespace
\multicolumn{5}{@{}l}{\textit{Central}} \\
Cz   & 0.00 &  0.00 & 1.000 & EEG \\
C3   & -0.35 &  0.00 & 0.937 & EEG \\
C4   & 0.35 &  0.00 & 0.937 & EEG \\
CP5  & -0.60 &  0.00 & 0.800 & EEG \\
CP1  & -0.60 &  0.00 & 0.800 & EEG \\
CP6  & 0.60 &  0.00 & 0.800 & EEG \\
CP2  & 0.60 &  0.00 & 0.800 & EEG \\
\addlinespace
\multicolumn{5}{@{}l}{\textit{Temporal}} \\
T7   & -0.95 &  0.00 & 0.312 & EEG \\
T8   & 0.95 &  0.00 & 0.312 & EEG \\
TP9  & -0.95 & -0.20 & 0.240 & EEG \\
TP10 & 0.95 & -0.20 & 0.240 & EEG \\
\addlinespace
\multicolumn{5}{@{}l}{\textit{Parietal}} \\
Pz   & 0.00 & -0.70 & 0.714 & EEG \\
P3   & -0.35 & -0.68 & 0.644 & EEG \\
P7   & -0.85 & -0.68 & 0.000 & EEG \\
P4   & 0.35 & -0.68 & 0.644 & EEG \\
P8   & 0.85 & -0.68 & 0.000 & EEG \\
\addlinespace
\multicolumn{5}{@{}l}{\textit{Occipital}} \\
O1   & -0.60 & -0.90 & 0.000 & EEG \\
Oz   & 0.00 & -0.90 & 0.436 & EEG \\
O2   & 0.60 & -0.90 & 0.000 & EEG \\
\addlinespace
FCz  & 0.00 & 0.35 & 0.937 & REF \\
Fpz  & 0.00 & 0.90 & 0.436 & GND \\
\bottomrule
\end{tabular}
\end{table}
